\section{问题重述}

\subsection{问题背景}
随着人工智能技术的快速发展，具身智能、智能交互和人机协同等应用逐渐从单纯的信息识别向更加深入的人类状态理解发展。在实际人机交互过程中，人的情感状态往往直接影响其语言表达、语音语调、面部表情以及肢体行为。因此，对人类情感进行准确识别与预测，是提高智能系统交互自然性和理解能力的重要基础。

与传统的单模态情感分析不同，真实环境中的情感表达通常由多种信息共同构成。文本模态主要反映语言内容及其语义信息，语音模态能够体现语调、节奏、音量等情感变化，视觉模态则可以通过面部表情和行为动作反映个体的外在情绪状态。不同模态之间具有较强的互补性，综合利用文本、语音和视觉信息，有助于更加全面地刻画人的情感状态。

然而，复杂场景下的多模态情感预测仍面临较大困难。一方面，文本、语音和视觉数据具有不同的采样频率、特征形式和时间尺度，不同模态的信息难以直接建立统一的时序对应关系；另一方面，在实际数据采集过程中，可能受到环境噪声、画面遮挡和语音转写异常等因素影响，使部分时段的模态信息缺失，从而降低传统多模态融合模型的稳定性。此外，现有情感预测模型通常更加关注最终预测结果，而对不同模态在决策过程中的作用以及关键情感证据缺乏直观、可量化的解释。

\subsection{问题重述}

\begin{enumerate}[label=\textbf{问题\arabic*:},leftmargin=*]
    \item 多模态特征提取与时序对齐,
从原始视频中分别提取文本、语音和视觉三模态情感特征，并建立统一的时序对应关系。保证各模态特征与原始样本、时间位置一一对应，为后续情感预测提供规范、可复现的输入。
    \item 模态局部缺失下的鲁棒情感预测,
针对文本、语音或视觉在部分连续时间段缺失的情况，建立能够动态利用剩余有效模态的鲁棒预测模型。同时分析缺失模态类型、缺失位置和缺失时长对情感极性分类与情感强度回归性能的影响。
    \item 可解释性情感预测,
建立既能输出情感极性和情感强度，又能解释预测依据的多模态模型。量化文本、语音、视觉三种模态的贡献，并定位对预测起关键作用的文本片段、语音时段和视觉关键帧。
\end{enumerate}
